\documentclass[../main.tex]{subfiles}
\begin{document}

\begin{center}
    \Large{\textbf{ABSTRACT}}\\
\end{center}
\vspace{1cm}

In recent years, Large Language Models (LLMs) have created new opportunities for building medical question-answering and clinical decision-support systems. However, medical applications require stronger guarantees than ordinary conversational software. A clinical assistant must reduce hallucination, preserve privacy, trace its answers to retrieved evidence, and support practical workflows such as laboratory-report interpretation, medication lookup, and follow-up reminders. A system that only forwards user questions to a generic cloud chatbot is therefore not sufficient for a safe and explainable healthcare prototype.

This thesis presents the updated implementation of a Smart Medical QA \& Consultation System that combines Retrieval-Augmented Generation (RAG), GraphRAG-oriented retrieval, Neo4j knowledge graph querying, and Agentic AI. The current source code is organized as a modular web application rather than a single chatbot script. The backend is implemented with FastAPI and exposes separate routes for health checking, Ollama proxying, graph retrieval, agent consultation, streaming responses, and medical-record analysis. The frontend is a browser-based interface containing pages for dashboard, chat, document upload, health news, and medication reminders.

The first core pillar of the system is a retrieval-grounded medical knowledge layer. Repository abstractions allow the agent and API endpoints to query a configured backend without being tightly coupled to a specific storage implementation. Neo4j and GraphRAG helper modules provide graph context retrieval, source metadata, evidence formatting, and answer synthesis. In parallel, the system integrates verified, custom-built knowledge graph artifacts containing a comprehensive medical ontology of clinical documents, text chunks, normalized entities, semantic mentions, and relations. These artifacts are produced by custom knowledge-graph modules for document chunking, entity extraction, relation extraction, predicate normalization, ontology utilities, and Neo4j ingestion. This design helps the system answer from retrieved context instead of relying only on the parametric memory of the language model.

The second pillar is a robust ReAct-style medical agent. The agent layer contains intent routing, retrieval planning, tool execution, ReAct parsing, parse-retry prompts, loop protection, forced finalization, confidence estimation, and streaming event generation. The available tools include graph retrieval and pill-image lookup, so the model can act on external evidence before writing the final response. The router also detects social turns and simple deterministic paths, allowing the system to avoid expensive agent execution when a full reasoning loop is unnecessary.

The third pillar is multi-format clinical-record processing. The \texttt{medical\_records} package supports PDF and Excel extraction, laboratory-value parsing, reference-range loading, on-form comparison, abnormal-result classification, medication-suggestion extraction, pill-image storage, and LLM-based advisory summaries. This turns the project from a pure chat application into a more complete clinical-support prototype that can analyze user-provided laboratory reports.

The fourth pillar is operational readiness and usability. Centralized settings manage Ollama, OpenRouter-compatible configuration, Neo4j, vector-store parameters, rate limiting, CORS, and agent behavior. The backend includes cache and Neo4j connection-pooling utilities, health/readiness endpoints, and tests for routes, services, repositories, settings, cache behavior, intent routing, ReAct logic, and evaluation utilities. The web UI supports local reminder scheduling through browser storage, making it possible for users to connect consultation output with medication-follow-up planning.

Overall, the updated system demonstrates a private, extensible, and explainable architecture for localized medical consultation. It does not claim to replace clinicians or provide certified diagnosis. Instead, it shows how local LLMs, graph retrieval, deterministic validation, document parsing, and user-facing workflows can be combined to reduce hallucination risk and improve traceability in a research-oriented CDSS prototype.

\vspace{0.5cm}
\noindent \textbf{Keywords:} Clinical Decision Support System (CDSS), Retrieval-Augmented Generation, GraphRAG, ReAct Agent, Neo4j, Ollama, Medical Record Analysis, Explainable AI.

\begin{flushright}
\begin{tabular}{@{}c@{}}
Vương Lê Bình Dương\\
\textit{(Signature and full name)}
\end{tabular}
\end{flushright}
\end{document}